\uuid{p4OW}
\titre{Dérivée en un point}
\auteur{bodin ; hanani ; mzari}
\organisation{exo7}
\licence{CC BY-NC-SA 4.0}
\datecreate{2020-01-20}
\niveau{L1}
\module{Analyse}
\chapitre{Dérivabilité des fonctions réelles}
\sousChapitre{Autre}
\difficulte{1}
\theme{Dérivées, dérivée en un point}
\exo7qcmid{lille-1.derivees.004}

\contenu{
\begin{question}
\qclassification{124.00}
\qtags{motcle=dérivée en un point}
Soit $\displaystyle f(x)=\sqrt[3]{(x-2)^2}$. Quelles sont les bonnes réponses ?
\begin{answers}
    \bad{$f$ est continue et dérivable en $2$.}
    \good{$f$ est continue et non dérivable en $2$.}
    \good{La tangente à $\mathcal{C}_f$ en $2$ est une droite verticale.}
    \bad{La tangente à $\mathcal{C}_f$ en $2$ est une droite horizontale.}
\end{answers}
\begin{explanations}
Les théorèmes généraux impliquent que $f$ est continue sur $\Rr$ et est dérivable sur $\Rr\setminus\{2\}$. Mais
$$\lim_{x\to 2^{\pm}}\frac{f(x)-f(2)}{x-2}=\lim_{x\to 2^{\pm}}\frac{1}{\sqrt[3]{x-2}}={\pm}\infty $$
Donc, $f$ n'est pas dérivable en $2$ et la tangente à $\mathcal{C}_f$ en $2$ est une droite verticale.
\end{explanations}
\end{question}
}
